\begin{solapartat}
\figura[0.6]{sol_fig1.png}
Només s'induirà corrent si es produeix una variació de flux magnètic a través de la superfície de l'espira \punts{0,2}\unskip.

La variació de flux es produirà en entrar i en sortir l'espira del camp magnètic. \punts{0,2}

D'acord amb la Llei de Lenz, el sentit del corrent serà tal que originaria un camp magnètic que s'oposarà a la variació de flux magnètic \punts{0,2}\unskip, de manera que:
\begin{itemize}[label=-]
  \item A l'entrar l'espira, el sentit de circulació del corrent serà antihorari \punts{0,2}\unskip.
  \item En sortir del camp magnètic el sentit del corrent serà horari \punts{0,2}\unskip.
\end{itemize}
\end{solapartat}

\begin{solapartat}
\punts{0,3} $\varepsilon = \left|\dfrac{d\varphi}{dt}\right| = \dfrac{B\,ds}{dt} = Ba\,\dfrac{dx}{dt} = Bav$ \ (en entrar i en sortir del camp)

\punts{0,3} $\varepsilon = 0,2\cdot 0,03\cdot 2 = 0,012\ \mathrm{V}$

\punts{0,4} $I = \dfrac{\varepsilon}{R} = \dfrac{0,012}{5} = 2,4\times 10^{-3}\ \mathrm{A}$ (en entrar i en sortir del camp)
\end{solapartat}
